% Section 3.2: 数据清洗与可执行校验

由 \texttt{convet\_KodCode\_to\_jsonl.py} 顺序扫描原始 Parquet，逐条同时执行格式与可执行性两层过滤，凑齐 \texttt{HEAD\_N=10000} 条目标样本。

\begin{itemize}
    \item \textbf{格式规范化}：要求 \texttt{question} / \texttt{solution} / \texttt{test} 三个字符串字段 \texttt{strip()} 后非空，且 \texttt{test\_info} 首元素为 dict 并同时含 \texttt{function\_declaration} / \texttt{function\_name} / \texttt{parameter\_list}（后续拼接评测 prompt 与少样本检索的最小必需字段）；
    \item \textbf{可执行校验}：脚本将 \texttt{solution} 与 \texttt{test} 写入临时目录，在隔离子进程中运行（带超时上限），全部断言通过且退出码为 $0$ 才记一次成功，其余（断言失败、导入错误、\texttt{SyntaxError}、超时）一律丢弃；
    \item \textbf{口径复用}：上述执行器 \texttt{test\_solution\_code} 同时被第四章 pass@$k$ 评测调用，避免清洗与评测口径漂移；
    \item \textbf{落盘}：通过校验的样本以固定字段序、\texttt{ensure\_ascii=False} 单行写入 \texttt{KodCode.jsonl}。
\end{itemize}

经此两层过滤，最终保留样本满足「官方 \texttt{solution} 能通过其自带 \texttt{test}」硬约束。该约束直接保证后续 3.4 节评测集上参考解答 pass@1 接近 $1$，作为模型侧实验的名义上界。
